\section{A cost model for agent collectives}
\label{sec:cost}

Collective algorithm selection is an optimisation problem, and changing the
objective changes the answer. This section states \ampi{}'s objective, works
out what it implies for the standard algorithms, and derives the one
constraint that has no MPI analogue.

\subsection{The model}

MPI's working model is Hockney's: a point-to-point message of $n$ bytes
costs $\alpha + n\beta$~\cite{hockney1994communication}. Computation is
omitted because it is negligible. For agents we need three terms, and the
one MPI omits dominates:

\begin{equation}
T \;=\; \gamma R \;+\; \alpha M \;+\; \beta N
\label{eq:cost}
\end{equation}

where $R$ is the number of \textbf{rounds} on the critical path (a round
being one agent turn: read the inputs, think, emit), $M$ the number of
messages, $N$ the number of tokens moved, and

\begin{itemize}
\item $\gamma$ is the cost of one agent turn: $10$--$100$ seconds and, at
2026 frontier-model prices, $10^{-2}$--$10^{0}$ currency units;
\item $\alpha$ is the per-message transport cost: $10^{-4}$--$10^{-2}$
seconds for a filesystem or HTTP hop;
\item $\beta$ is the per-token transport cost, effectively $10^{-7}$
seconds.
\end{itemize}

So $\gamma/\alpha \in [10^3, 10^6]$ and $\gamma/(\beta \cdot 10^4) \approx
10^5$. Communication is free; turns are everything. \Cref{eq:cost} is
therefore, to within a rounding error,

\begin{equation}
T \approx \gamma R.
\label{eq:cost-simple}
\end{equation}

Two secondary objectives remain, and both are budgets rather than costs:
total tokens ingested (which is money, and also context) and total turns
(which is money). We report all three.

\subsection{What \texorpdfstring{$T \approx \gamma R$}{T = gamma R} does to
algorithm selection}

\Cref{tab:rounds} gives round counts for the standard algorithms.
Under \Cref{eq:cost-simple} the column that matters is the first, and the
consequences are stark.

\begin{table}[t]
\centering
\small
\caption{Rounds on the critical path and per-rank ingest, $p$ ranks,
contribution size $s$ tokens. Under $T \approx \gamma R$ the leftmost
column decides wall-clock and the rightmost decides feasibility.}
\label{tab:rounds}
\begin{tabular}{@{}llll@{}}
\toprule
\textbf{Collective} & \textbf{Algorithm} & \textbf{Rounds} & \textbf{Peak ingest}\\
\midrule
Barrier & linear & $2$ & $0$\\
        & dissemination & $\lceil\log_2 p\rceil$ & $0$\\
\midrule
Bcast & flat & $1$ & $s$\\
      & binomial & $\lceil\log_2 p\rceil$ & $s$\\
      & chain & $p-1$ & $s$\\
\midrule
Gather & flat & $1$ & $(p-1)s$ at root\\
       & binomial & $\lceil\log_2 p\rceil$ & $(p/2)s$ at root\\
\midrule
Allgather & ring & $p-1$ & $(p-1)s$\\
          & rec.\ doubling & $\lceil\log_2 p\rceil$ & $(p-1)s$\\
          & Bruck & $\lceil\log_2 p\rceil$ & $(p-1)s$\\
\midrule
Reduce & flat & $1$ & $(p-1)s$ at root\\
       & $k$-ary tree & $\lceil\log_k p\rceil$ & $(k-1)\lceil\log_k p\rceil\,m$\\
\midrule
Allreduce & ring & $2(p-1)$ & $s$\\
          & rec.\ doubling & $\lceil\log_2 p\rceil$ & $s\log_2 p$\\
          & reduce+bcast & $\lceil\log_k p\rceil{+}\lceil\log_2 p\rceil$ & tree\\
\midrule
Scan & chain & $p-1$ & $s$\\
     & Hillis--Steele & $\lceil\log_2 p\rceil$ & $s\log_2 p$\\
\midrule
Alltoall & pairwise & $p-1$ & $(p-1)s$\\
\bottomrule
\end{tabular}
\end{table}

\paragraph{The bandwidth-optimal crossover disappears.}
In MPI, ring allreduce is preferred at large message sizes because it is
bandwidth optimal: it moves $2\frac{p-1}{p}n$ bytes against recursive
doubling's $n\log_2 p$~\cite{rabenseifner2004optimization,Patarasuk2009}.
Under \Cref{eq:cost-simple} that advantage is worth nothing, while its
$2(p-1)$ rounds against $\log_2 p$ costs everything: at $p=64$, 126 turns
versus 6, or roughly twenty minutes versus one. There is no message size at
which the ring wins, and \Cref{sec:eval-micro} confirms the crossover does
not appear.

\paragraph{Flat algorithms become attractive again.}
Symmetrically, a flat broadcast---the root sends $p-1$ messages---costs one
round, against $\log_2 p$ for a tree. In MPI the flat form is dismissed
because the root serialises $p-1$ sends, each costing $\alpha$; here those
sends are file writes and the root's cost is one turn regardless. The tree
is preferable only when the root's *sending* is itself a turn (a refracting
broadcast) or when message fan-out is rate-limited. Our decision function
selects flat for $p \le 4$ and binomial above, and the honest statement is
that the crossover is set by provider concurrency limits rather than by the
cost model.

\paragraph{Scan is the highest-leverage collective.}
A harness with a carried dependency---each worker needs something every
earlier worker decided---is usually written as a chain, at $p-1$ turns. As a
prefix scan it is $\lceil\log_2 p\rceil$. At $p=12$ that is 4 rounds instead
of 11; at $p=64$, 6 instead of 63. Nothing else in the catalogue offers that
factor, and it applies to a structure (shared conventions established
incrementally) that appears in almost every non-trivial agent workload.

\subsection{Feasibility: the constraint MPI does not have}
\label{sec:cost-cumulative}

Every MPI collective is feasible. Under M1 (\Cref{sec:m1}) that stops being
true, and feasibility must be checked before cost.

Let $C$ be a rank's usable context budget---its window less the reserve it
needs for its own reasoning and output. Then:

\begin{observation}[Allgather is intrinsically limited]
\op{Allgather} requires every rank to ingest $(p-1)s$ tokens, by definition
and not by algorithm. It is therefore infeasible for
$p > 1 + C/s$, and no choice of schedule helps. With $C = 80{,}000$ and
$s = 2{,}000$, that is $p = 41$.
\end{observation}

This is worth stating because ``let every agent see every draft'' is the
first thing a harness author reaches for, and it has a hard ceiling that
arrives sooner than intuition suggests. \op{Allreduce} with a contracting
operator is the scalable replacement, and the runtime says so by name when
it predicts the overflow.

The reduction case is more subtle, and it is where we made a mistake worth
recording. A $k$-ary reduction tree over $p$ ranks has depth
$d=\lceil\log_k p\rceil$, and each interior node receives $k-1$
contributions \emph{per round}. In MPI the memory cost is one round's worth,
$(k-1)m$, because the buffers are reused; depth costs only latency. We
initially planned trees against that figure. It is wrong:

\begin{observation}[Tree depth costs context]
An agent retains everything it reads, so the root of a depth-$d$ $k$-ary
tree ingests
\begin{equation}
I_{\mathrm{root}} = (k-1)\,d\,m
= (k-1)\lceil\log_k p\rceil\, m
\label{eq:cumulative}
\end{equation}
tokens over the life of the reduction, where $m$ is the operator's output
bound. Planning against the per-round figure $(k-1)m$ yields schedules that
look feasible, run two rounds, and exhaust the root.
\end{observation}

\Cref{eq:cumulative} is not monotone in $k$, which makes the planning
problem genuinely non-trivial: increasing $k$ raises the per-round fan-in
but lowers the depth. The planner therefore searches $k \in [2,p]$ for the
schedule minimising $d$ subject to $I_{\mathrm{root}} \le C$, breaking ties
on $I_{\mathrm{root}}$. A worked example: $p=64$, $m=1000$, $C=10{,}000$. A
degree-10 tree finishes in 2 rounds but costs the root
$9 \times 2 \times 1000 = 18{,}000$ tokens---over budget. Degree 4 costs
$3 \times 3 \times 1000 = 9{,}000$ in 3 rounds, and is selected. The
protocol trades one round for feasibility, which is exactly the trade a
harness author cannot make by hand because they cannot see the constant.

Finally, the operator's contraction is what makes any of this possible:

\begin{observation}[Non-contracting operators forbid deep trees]
If the operator's output is no smaller than its inputs, $m$ grows with
depth, $I_{\mathrm{root}}$ grows super-linearly, and only a flat reduction
is safe. A flat reduction costs the root $(p-1)s$, so the whole reduction is
limited to $p \le 1 + C/s$ ranks.
\end{observation}

This is why \ampi{} requires an operator to declare an output bound rather
than inferring one: the declaration is what turns an unbounded reduction
into a bounded one, and the runtime cannot discover it by observation
without first running the thing it is trying to plan.

\subsection{What the model does not capture}

Three caveats, stated because we rely on the model in \Cref{sec:eval}.

First, $\gamma$ is not constant. An agent turn's duration grows with input
length and output length, so a rank deep in a reduction tree, holding more
context, is slower. The model treats $\gamma$ as a constant and therefore
under-predicts the cost of context-heavy schedules---in the direction that
makes our capacity-aware planning look \emph{less} valuable than it is.

Second, turn durations are heavy-tailed. A collective's cost is the
\emph{maximum} over participants, not the mean, so $\gamma$ in
\Cref{eq:cost} should be read as a high quantile and the gap between mean
and max widens with $p$. This is the standard tail-at-scale
effect~\cite{Dean2013}, and it is the argument for \op{Waitsome} over
\op{Waitall} wherever the harness can proceed with partial results.

Third, the model says nothing about quality. A schedule that is cheaper may
produce a worse answer---a deeper reduction tree summarises more
aggressively and loses more---so cost and quality must be reported together.
\Cref{sec:eval} does.
